\section{Introduction}

\subsection{Problem}
Network monitoring and intrusion detection rarely work on single packets. They
work on \emph{flows}: all the packets that share the same source address,
destination address, source port, destination port and protocol. A flow is
summarized by a record of features such as its packet and byte counts, the
statistics of its packet sizes and inter-arrival times, the Transmission Control
Protocol (TCP) flags it carried, and how bursty it was. Flow records are the
input of flow exporters such as NetFlow~\cite{netflow}, of feature extractors
such as CICFlowMeter~\cite{cicflowmeter}, and of the machine-learning detectors
trained on data sets such as CIC-IDS2017~\cite{cicids}.

The problem we solve is the step from packets to flow records. Given a batch of
captured packets in time order, we must group the packets by flow and compute
the features of every flow. A backbone link carries hundreds of thousands of
packets per second, so a capture of a few minutes already holds tens of
millions of packets and millions of flows.

\subsection{Why it is important and why it is hard}
Feature extraction is the first stage of every flow-based detector, and it
must keep up with the traffic. If it is slow, the detector sees the traffic
late or sees only a sample of it. The computation is simple per packet, but
three properties make it expensive and hard to parallelize.

\emph{Irregular memory access.} Every packet updates the record of its flow,
and consecutive packets belong to unrelated flows. With millions of flows the
records do not fit in the processor caches, so most updates are cache misses.

\emph{Skew.} Traffic is heavy-tailed: a few ``elephant'' flows carry a large
share of the packets while most flows are ``mice'' with one or two
packets~\cite{elephants}. When many packets belong to the same flow, threads
that update one shared table contend for the same record~\cite{cieslewicz}, and
schemes that give each flow to one worker become unbalanced.

\emph{Order dependence.} Inter-arrival times and bursts depend on the order of
the packets inside a flow. A parallel method cannot simply add contributions
in any order; it has to keep or rebuild the time order of each flow.

\subsection{Background}
Grouping and aggregation is a classic database operation. On multicore
processors the two main families are a shared hash table and partitioning the
input so that each thread owns a disjoint part of the
keys~\cite{balkesen,cieslewicz}. On a Graphics Processing Unit (GPU), grouping
is usually built from data-parallel primitives: lock-free hash tables filled
with compare-and-swap instructions~\cite{alcantara}, radix sort, prefix sums
and segmented reductions, available for example in the Thrust
library~\cite{thrust}. We use three programming models from the course:
OpenMP~\cite{openmp} for shared-memory threads, the Compute Unified Device
Architecture (CUDA)~\cite{cuda} for the GPU, and the Message Passing Interface
(MPI) through its Python binding MPI4Py~\cite{mpi4py} for distributing batches
of packet capture (PCAP) files among processes.

\subsection{Contributions}
\begin{itemize}
\item A packet-to-flow pipeline that computes \nfeat{} features per flow,
including order-dependent ones, with a sequential baseline, two OpenMP
strategies, a CUDA pipeline with three reduction kernels, and an MPI4Py
driver for batches of capture files.
\item A formulation of the per-flow record as an associative merge, which lets
threads process time-contiguous chunks independently and still obtain exact
inter-arrival and burst features.
\item An exactness guarantee: every parallel version returns the same integer
accumulators as the baseline on every input tested (\nchecks{} automated
checks), and derived features equal to within floating-point rounding.
\item An experimental study on synthetic traffic with controlled skew and on a
real backbone trace, covering runtime, speedup, efficiency, strong and weak
scaling, GPU launch configurations, transfer cost, contention and
communication overhead. The main finding is that no single strategy wins:
the best choice depends on the number of flows and on the skew.
\end{itemize}
